\documentclass[reqno]{amsart}
\usepackage[T1]{fontenc}
\usepackage{amsmath,amssymb,amsthm}
\usepackage{xr-hyper}
\usepackage[hidelinks]{hyperref}
\externaldocument[ec-]{20261124T223000--elliptic-curves-for-the-proof__workbook}
\externaldocument[cm-]{20261124T223100--complex-multiplication-and-inert-reduction__workbook}
\externaldocument[ip-]{20261124T214600--the-isotropic-plane__workbook}

% Chapter number; \ref{chap} gives it.
\newcounter{chap}
% Chapter numbers, generated from 20261125T012000--grh-workbook__guide.org
% by scripts/make-chapters.py.  Do not edit: change the order of the
% index and rerun the script (./preview-site.sh and CI do this).
%
% Usage in a sheet, after \newcounter{chap}:
%   % Chapter numbers, generated from 20261125T012000--grh-workbook__guide.org
% by scripts/make-chapters.py.  Do not edit: change the order of the
% index and rerun the script (./preview-site.sh and CI do this).
%
% Usage in a sheet, after \newcounter{chap}:
%   % Chapter numbers, generated from 20261125T012000--grh-workbook__guide.org
% by scripts/make-chapters.py.  Do not edit: change the order of the
% index and rerun the script (./preview-site.sh and CI do this).
%
% Usage in a sheet, after \newcounter{chap}:
%   \input{chapters}\chapnumber{<sheet ID>}
\makeatletter
\providecommand\chapters@set[2]{\@namedef{chapters@#1}{#2}}
\chapters@set{20261125T001000}{1}% zeros-to-ramp-sums
\chapters@set{20261125T001100}{2}% inert-primes-and-estimate-b
\chapters@set{20261125T005000}{3}% galois-fourier-and-averaging
\chapters@set{20261125T005500}{4}% columns-from-a-sphere
\chapters@set{20261125T000000}{5}% slope-method
\chapters@set{20261126T030100}{6}% a-twist-for-each-prime
\chapters@set{20261125T003000}{7}% the-chain-and-the-square-root
\chapters@set{20261125T010000}{8}% rank-lemma-and-zero-estimate
\chapters@set{20261125T010500}{9}% poincare-series-and-near-orthogonality
\chapters@set{20261124T223000}{10}% elliptic-curves-for-the-proof
\chapters@set{20261125T009500}{11}% fock-model-and-theta-intertwiners
\chapters@set{20261124T223100}{12}% complex-multiplication-and-inert-reduction
\chapters@set{20261125T008000}{13}% one-cm-determinant
\chapters@set{20261124T214500}{14}% supersingular-theta-divisibility
\chapters@set{20261124T214600}{15}% the-isotropic-plane
\chapters@set{20261125T013700}{16}% pairings-on-a-formal-group-of-height-two
\chapters@set{20261125T013500}{17}% hecke-nondegeneracy-on-the-tree
\chapters@set{20261126T030200}{18}% the-bound-is-attained
\chapters@set{20261125T009000}{19}% kronecker-limit-and-chowla-selberg
\chapters@set{20261125T002000}{20}% definite-theta-lifts
\chapters@set{20261125T015500}{21}% siegel-weil-in-the-smallest-case
\chapters@set{20261125T002100}{22}% definite-theta-lifts-and-norms
\chapters@set{20261125T016500}{23}% hecke-operators-and-theta-lifts
\chapters@set{20261125T016700}{24}% divisibility-along-the-chain
\chapters@set{20261125T004000}{25}% denominators-and-sha
\chapters@set{20261125T004100}{26}% tate-twists-and-sha
\chapters@set{20261126T020430}{27}% the-minus-part-by-gauss-sums
\chapters@set{20261125T014100}{28}% congruence-modules
\chapters@set{20261125T014500}{29}% eisenstein-congruences-and-selmer-classes
\chapters@set{20261125T014700}{30}% pseudocharacters-and-the-length-bound
\chapters@set{20261125T011000}{31}% galois-representations-black-box
\chapters@set{20261126T030300}{32}% the-data
\chapters@set{20261126T020210}{33}% bessel-poincare-series-and-their-gram-operators
\chapters@set{20261126T030400}{34}% the-seesaw-identity
\chapters@set{20261126T030500}{35}% divisibility-of-the-composite-lift
\chapters@set{20261126T030600}{36}% theorem-a-short-forms
\chapters@set{20261126T030700}{37}% the-determinant-of-a-chain
\chapters@set{20261126T030800}{38}% the-cost-of-one-lift
\chapters@set{20261126T030900}{39}% denominators-of-projections-to-lifts
\chapters@set{20261126T031000}{40}% twists-with-few-congruences
\chapters@set{20261126T031100}{41}% theorem-c-and-the-main-inequality
\providecommand\chapnumber[1]{%
  \@ifundefined{chapters@#1}%
    {\PackageError{chapters}{Sheet #1 is not listed in the index}%
      {List it in 20261125T012000--grh-workbook__guide.org\MessageBreak
       and run scripts/make-chapters.py.}}%
    {\edef\thechap{\@nameuse{chapters@#1}}}}
\makeatother
\chapnumber{<sheet ID>}
\makeatletter
\providecommand\chapters@set[2]{\@namedef{chapters@#1}{#2}}
\chapters@set{20261125T001000}{1}% zeros-to-ramp-sums
\chapters@set{20261125T001100}{2}% inert-primes-and-estimate-b
\chapters@set{20261125T005000}{3}% galois-fourier-and-averaging
\chapters@set{20261125T005500}{4}% columns-from-a-sphere
\chapters@set{20261125T000000}{5}% slope-method
\chapters@set{20261126T030100}{6}% a-twist-for-each-prime
\chapters@set{20261125T003000}{7}% the-chain-and-the-square-root
\chapters@set{20261125T010000}{8}% rank-lemma-and-zero-estimate
\chapters@set{20261125T010500}{9}% poincare-series-and-near-orthogonality
\chapters@set{20261124T223000}{10}% elliptic-curves-for-the-proof
\chapters@set{20261125T009500}{11}% fock-model-and-theta-intertwiners
\chapters@set{20261124T223100}{12}% complex-multiplication-and-inert-reduction
\chapters@set{20261125T008000}{13}% one-cm-determinant
\chapters@set{20261124T214500}{14}% supersingular-theta-divisibility
\chapters@set{20261124T214600}{15}% the-isotropic-plane
\chapters@set{20261125T013700}{16}% pairings-on-a-formal-group-of-height-two
\chapters@set{20261125T013500}{17}% hecke-nondegeneracy-on-the-tree
\chapters@set{20261126T030200}{18}% the-bound-is-attained
\chapters@set{20261125T009000}{19}% kronecker-limit-and-chowla-selberg
\chapters@set{20261125T002000}{20}% definite-theta-lifts
\chapters@set{20261125T015500}{21}% siegel-weil-in-the-smallest-case
\chapters@set{20261125T002100}{22}% definite-theta-lifts-and-norms
\chapters@set{20261125T016500}{23}% hecke-operators-and-theta-lifts
\chapters@set{20261125T016700}{24}% divisibility-along-the-chain
\chapters@set{20261125T004000}{25}% denominators-and-sha
\chapters@set{20261125T004100}{26}% tate-twists-and-sha
\chapters@set{20261126T020430}{27}% the-minus-part-by-gauss-sums
\chapters@set{20261125T014100}{28}% congruence-modules
\chapters@set{20261125T014500}{29}% eisenstein-congruences-and-selmer-classes
\chapters@set{20261125T014700}{30}% pseudocharacters-and-the-length-bound
\chapters@set{20261125T011000}{31}% galois-representations-black-box
\chapters@set{20261126T030300}{32}% the-data
\chapters@set{20261126T020210}{33}% bessel-poincare-series-and-their-gram-operators
\chapters@set{20261126T030400}{34}% the-seesaw-identity
\chapters@set{20261126T030500}{35}% divisibility-of-the-composite-lift
\chapters@set{20261126T030600}{36}% theorem-a-short-forms
\chapters@set{20261126T030700}{37}% the-determinant-of-a-chain
\chapters@set{20261126T030800}{38}% the-cost-of-one-lift
\chapters@set{20261126T030900}{39}% denominators-of-projections-to-lifts
\chapters@set{20261126T031000}{40}% twists-with-few-congruences
\chapters@set{20261126T031100}{41}% theorem-c-and-the-main-inequality
\providecommand\chapnumber[1]{%
  \@ifundefined{chapters@#1}%
    {\PackageError{chapters}{Sheet #1 is not listed in the index}%
      {List it in 20261125T012000--grh-workbook__guide.org\MessageBreak
       and run scripts/make-chapters.py.}}%
    {\edef\thechap{\@nameuse{chapters@#1}}}}
\makeatother
\chapnumber{<sheet ID>}
\makeatletter
\providecommand\chapters@set[2]{\@namedef{chapters@#1}{#2}}
\chapters@set{20261125T001000}{1}% zeros-to-ramp-sums
\chapters@set{20261125T001100}{2}% inert-primes-and-estimate-b
\chapters@set{20261125T005000}{3}% galois-fourier-and-averaging
\chapters@set{20261125T005500}{4}% columns-from-a-sphere
\chapters@set{20261125T000000}{5}% slope-method
\chapters@set{20261126T030100}{6}% a-twist-for-each-prime
\chapters@set{20261125T003000}{7}% the-chain-and-the-square-root
\chapters@set{20261125T010000}{8}% rank-lemma-and-zero-estimate
\chapters@set{20261125T010500}{9}% poincare-series-and-near-orthogonality
\chapters@set{20261124T223000}{10}% elliptic-curves-for-the-proof
\chapters@set{20261125T009500}{11}% fock-model-and-theta-intertwiners
\chapters@set{20261124T223100}{12}% complex-multiplication-and-inert-reduction
\chapters@set{20261125T008000}{13}% one-cm-determinant
\chapters@set{20261124T214500}{14}% supersingular-theta-divisibility
\chapters@set{20261124T214600}{15}% the-isotropic-plane
\chapters@set{20261125T013700}{16}% pairings-on-a-formal-group-of-height-two
\chapters@set{20261125T013500}{17}% hecke-nondegeneracy-on-the-tree
\chapters@set{20261126T030200}{18}% the-bound-is-attained
\chapters@set{20261125T009000}{19}% kronecker-limit-and-chowla-selberg
\chapters@set{20261125T002000}{20}% definite-theta-lifts
\chapters@set{20261125T015500}{21}% siegel-weil-in-the-smallest-case
\chapters@set{20261125T002100}{22}% definite-theta-lifts-and-norms
\chapters@set{20261125T016500}{23}% hecke-operators-and-theta-lifts
\chapters@set{20261125T016700}{24}% divisibility-along-the-chain
\chapters@set{20261125T004000}{25}% denominators-and-sha
\chapters@set{20261125T004100}{26}% tate-twists-and-sha
\chapters@set{20261126T020430}{27}% the-minus-part-by-gauss-sums
\chapters@set{20261125T014100}{28}% congruence-modules
\chapters@set{20261125T014500}{29}% eisenstein-congruences-and-selmer-classes
\chapters@set{20261125T014700}{30}% pseudocharacters-and-the-length-bound
\chapters@set{20261125T011000}{31}% galois-representations-black-box
\chapters@set{20261126T030300}{32}% the-data
\chapters@set{20261126T020210}{33}% bessel-poincare-series-and-their-gram-operators
\chapters@set{20261126T030400}{34}% the-seesaw-identity
\chapters@set{20261126T030500}{35}% divisibility-of-the-composite-lift
\chapters@set{20261126T030600}{36}% theorem-a-short-forms
\chapters@set{20261126T030700}{37}% the-determinant-of-a-chain
\chapters@set{20261126T030800}{38}% the-cost-of-one-lift
\chapters@set{20261126T030900}{39}% denominators-of-projections-to-lifts
\chapters@set{20261126T031000}{40}% twists-with-few-congruences
\chapters@set{20261126T031100}{41}% theorem-c-and-the-main-inequality
\providecommand\chapnumber[1]{%
  \@ifundefined{chapters@#1}%
    {\PackageError{chapters}{Sheet #1 is not listed in the index}%
      {List it in 20261125T012000--grh-workbook__guide.org\MessageBreak
       and run scripts/make-chapters.py.}}%
    {\edef\thechap{\@nameuse{chapters@#1}}}}
\makeatother
\chapnumber{20261124T214500}
\renewcommand{\thesection}{\thechap.\arabic{section}}

\theoremstyle{definition}
\newtheorem{problem}{Problem}[section]
\theoremstyle{plain}
\newtheorem{theorem}{Theorem}[section]

\newcommand{\Z}{\mathbb{Z}}
\newcommand{\Q}{\mathbb{Q}}
\newcommand{\C}{\mathbb{C}}
\newcommand{\F}{\mathbb{F}}
\newcommand{\ii}{\mathrm{i}}
\newcommand{\Oc}{\mathcal{O}}
\newcommand{\sll}{\operatorname{sl}}

\setlength{\emergencystretch}{2em}

\title{The invariant derivation at a supersingular prime}
\author{}
\date{November 24, 2026}

\begin{document}
\maketitle
\refstepcounter{chap}\label{chap}

\paragraph{Prerequisites.}
Formal groups of elliptic curves, their height, and supersingular reduction \cite[Ch.~IV and \S\S V.3--V.4]{Si1}.  The Weierstrass functions \(\wp\) and \(\sigma\) \cite[\S I.5]{Si2}.  Legendre's formula \(v_p(n!) = \sum_{i \ge 1} \lfloor n/p^i \rfloor\), Kummer's theorem that \(v_p\binom{m+n}{m}\) is the number of carries in the base-\(p\) addition of \(m\) and \(n\), and the Lagrange inversion formula.  From Chapter~\ref{ec-chap}, \emph{Complex tori, theta functions and periods}, we use the lattice \(\varpi\Z[\ii]\) of Problem~\ref{ec-pr:varpi}.  From Chapter~\ref{cm-chap}, \emph{Complex multiplication and reduction at an inert prime}: formal groups, their height and the invariant derivation (Section~\ref{cm-sec:formal}), and Problem~\ref{cm-pr:lambda}; Section~\ref{sec:Dp} proves, for the lemniscate, the congruence \(D^p \equiv \lambda D \pmod p\) stated there.  Chapter~\ref{ip-chap}, \emph{The isotropic plane}, continues this one.

\section{Gaussian moments}\label{sec:moments}

Let \(\zeta = e^{2\pi\ii/8}\), and for integers \(a, w \ge 0\) consider the moments
\[
  m(a,w) := \sum_{x \in \Z[\ii]^2} (x_1 + \zeta x_2)^a\, (x_1 - \zeta x_2)^w\, e^{-\pi(\lvert x_1 \rvert^2 + \lvert x_2 \rvert^2)} .
\]
After a normalization they are rational numbers.  How divisible are they by a prime \(p \equiv 3 \pmod 4\), that is, by a prime that is inert in \(\Q(\ii)\)?

Let \(\sll\) be the lemniscate sine, the meromorphic function with \(\sll(0) = 0\), \(\sll'(0) = 1\) and \(\sll'^2 = 1 - \sll^4\), and let \(\varpi := 2\int_0^1 (1 - t^4)^{-1/2} \, d t\).  Let \(\wp\) and \(\sigma\) be the Weierstrass functions attached to \(\Lambda = \varpi\Z[\ii]\), whose invariants are \(g_2 = 4\) and \(g_3 = 0\) (Problem~\ref{ec-pr:varpi}).  Let
\[
  F := \sigma^2 \wp .
\]

\begin{problem}\label{pr:wp-sl}
\(\wp = \sll^{-2}\).
\end{problem}

\begin{problem}\label{pr:F}
\(F\) is entire and even, \(F(\ii u) = F(u)\), and \(F(u) = 1 + \tfrac16 u^4 + \tfrac{17}{2520} u^8 + O(u^{12})\).
\end{problem}

\begin{problem}\label{pr:gauss}
Let \(\vartheta_0 := \sum_{\alpha \in \Z[\ii]} e^{-\pi \lvert \alpha \rvert^2}\).  For \(z \in \C\),
\[
  \sum_{\alpha \in \Z[\ii]} e^{-\pi\lvert\alpha\rvert^2 + z\alpha} = \vartheta_0\, F\Bigl(\frac{\varpi z}{2\pi}\Bigr).
\]
\end{problem}

Let \(\Theta(Z_1, Z_2) := F(Z_1) F(Z_2)\) on \(\C^2\), let \(e_\pm := (1, \pm\zeta)\), and let
\[
  t(a,w) := \partial_s^a \partial_t^w \, \Theta(s e_+ + t e_-) \big|_{s = t = 0} ,
  \qquad
  \Theta(s e_+ + t e_-) = F(s+t)\, F\bigl(\zeta(s-t)\bigr).
\]

\begin{problem}\label{pr:moments}
\(m(a,w) = \vartheta_0^2\, (\varpi/2\pi)^{a+w}\, t(a,w)\).
\end{problem}

\begin{problem}\label{pr:t-sym}
The numbers \(t(a,w)\) are rational, \(t(a,w) = t(w,a)\), and \(t(a,w) = 0\) unless \(a + 5w \equiv 0 \pmod 8\).
\end{problem}

For instance \(t(3,1) = 8\), \(t(4,4) = 448\) and \(t(8,0) = -576\).  The table lists \(v_{11}(t(a,w))\) for \(w \le 2\) and the first nine values of \(a \ge 8\) allowed by Problem~\ref{pr:t-sym}.
\begin{center}
\begin{tabular}{c|ccccccccc}
\(w = 0\), \(a =\) & 8 & 16 & 24 & 32 & 40 & 48 & 56 & 64 & 72 \\
\(v_{11}\) & 0 & 1 & 2 & 2 & 3 & 4 & 5 & 5 & 6 \\
\hline
\(w = 1\), \(a =\) & 11 & 19 & 27 & 35 & 43 & 51 & 59 & 67 & 75 \\
\(v_{11}\) & 0 & 0 & 1 & 2 & 2 & 3 & 4 & 5 & 5 \\
\hline
\(w = 2\), \(a =\) & 14 & 22 & 30 & 38 & 46 & 54 & 62 & 70 & 78 \\
\(v_{11}\) & 0 & 0 & 0 & 1 & 2 & 2 & 3 & 4 & 5
\end{tabular}
\end{center}
In this range \(v_{11}(t(a,w)) = \max(\lfloor a/11 \rfloor - w, 0)\).  Each derivative in the second direction seems to cost one factor of \(11\).  By contrast, \(t(8,0) = -576\) is prime to \(7\).  This chapter treats one variable: what \(p\) divides the derivatives of a function on the lemniscate, and why \(F\) alone does not gain.  The sequel, \emph{The isotropic plane}, proves a lower bound of the shape suggested by the \(11\)-adic valuations in the table at every prime \(p \equiv 3 \pmod 8\).

\section{The lemniscate and its formal group}\label{sec:lemniscate}

Let \(E\) be the smooth projective model of \(y^2 = 1 - X^4\), with origin \((0,1)\).  It is a curve of genus one with good reduction at every odd prime.  The map \(u \mapsto (\sll u, \sll' u)\) identifies \(\C/L\) with \(E(\C)\), where \(L \subset \Lambda\) is the period lattice of \(\sll\).  The function \(X\) is a parameter at the origin, \(du = dX/y\) is an invariant differential, and the addition law is
\[
  \sll(u + v) = \frac{\sll u\, \sll' v + \sll v\, \sll' u}{1 + \sll^2 u\, \sll^2 v} .
\]
Let \(D := d/du = y \, d/dX\) be the invariant derivation, and let \(\ell(X) := \int_0^X (1 - v^4)^{-1/2} \, dv\) be the logarithm, so that \(\sll(\ell(X)) = X\).

\begin{problem}\label{pr:sl-series}
\(\sll u = u - \tfrac1{10}u^5 + \tfrac1{120}u^9 - \tfrac{11}{15600}u^{13} + \dotsb\).  Only exponents \(\equiv 1 \pmod 4\) occur, \(\sll(\ii u) = \ii\, \sll u\), and the lattice \(L\) satisfies \(\ii L = L\).
\end{problem}

\begin{problem}\label{pr:fg}
The formal group law \(X \oplus X'\), defined by \(\sll u \oplus \sll v = \sll(u + v)\), has coefficients in \(\Z[\tfrac12]\), and after extension of scalars to \(\Z[\tfrac12,\ii]\) it has the endomorphism \(X \mapsto \ii X\).
\end{problem}

\begin{problem}\label{pr:D}
\(D\) maps \(\Z[\tfrac12][[X]]\) into itself, and \((D^n f)(0) = n!\, [u^n]\, f(\sll u)\) for \(f \in \Q[[X]]\).
\end{problem}

Fix an odd prime \(p\) and an embedding \(\Q(\zeta) \hookrightarrow \overline{\Q}_p\), and write \(\zeta\) and \(\ii = \zeta^2\) also for their images.  Let \(\Oc := \Z_p[\ii] \subseteq \overline{\Q}_p\), and let \(v_p\) be the \(p\)-adic valuation normalized by \(v_p(p) = 1\).  For \(p \equiv 3 \pmod 4\), \(\Oc\) is the ring of integers of the unramified quadratic extension of \(\Q_p\), and it contains \(\zeta\), since \(8\) divides \(p^2 - 1\).  We write \(\hat E\) for the formal group of Problem~\ref{pr:fg} over \(\Oc\).

\begin{problem}\label{pr:O-action}
For each \(\alpha \in \Oc\) the series \([\alpha](X) := \sll(\alpha \ell(X))\) lies in \(\Oc[[X]]\).  It is the unique endomorphism of \(\hat E\) with \([\alpha](X) = \alpha X + O(X^2)\).
\end{problem}

When \(p \equiv 3 \pmod 4\), this action of \(\Oc\) makes \(\hat E\) a formal \(\Oc\)-module of the kind constructed by Lubin and Tate \cite{LT}; we use only Problem~\ref{pr:O-action}.  For these primes \(E\) is supersingular and \(\hat E\) has height two \cite[Thm.~V.3.1]{Si1}.  Indeed \(E\) is isogenous to \(y^2 = x^3 - x\) by \((X, y) \mapsto (X^{-2}, -yX^{-3})\), and that curve is supersingular exactly when \(p \equiv 3 \pmod 4\) \cite[\S V.4]{Si1}.  Section~\ref{sec:Dp} shows what height two does to \(D\).

\section{The \texorpdfstring{\(p\)}{p}th power of the invariant derivation}\label{sec:Dp}

Let \(f\) be a function on \(E\), regular at the origin, with expansion in \(\Oc[[X]]\).  Its derivatives \((D^n f)(0)\) are integral by Problem~\ref{pr:D}.  We show that they are divisible by \(p^{\lfloor n/p \rfloor}\) when \(p \equiv 3 \pmod 4\).  The smallest case is \(p = 3\): \(DX = y\), \(D^2X = -2X^3\) and \(D^3X = -6X^2y\).

For a series \(f\), let \(\tau_Y f(X) := f(X \oplus Y)\), a series in two variables.

\begin{problem}\label{pr:Dp-derivation}
The operator induced by \(D^p\) on \((\Oc/p\Oc)[[X]]\) is a derivation.
\end{problem}

\begin{problem}\label{pr:invariant}
Let \(k\) be a field of characteristic \(p\), and write \(D\) and \(\tau_Y\) for the operators obtained by reducing the coefficients of the invariant derivation and the formal group law from \(\Z[\tfrac12]\) to \(k\).  Let \(\delta\) be a continuous \(k\)-linear derivation of \(k[[X]]\), acting on \(k[[X,Y]]\) with \(Y\) held fixed.  If \(\delta \tau_Y = \tau_Y \delta\), then \(\delta = (\delta X)(0)\, D\).  The operator \(D\) itself commutes with \(\tau_Y\).
\end{problem}

Thus \(D^p \equiv \lambda_p D \pmod p\), with \(\lambda_p \equiv (D^p X)(0) = p!\, [u^p] \sll u\).

\begin{problem}\label{pr:Dp}
If \(p \equiv 3 \pmod 4\), then \(\lambda_p \equiv 0\), so that \(D^p\, \Oc[[X]] \subseteq p\, \Oc[[X]]\).
\end{problem}

\begin{problem}\label{pr:first-gain}
Let \(p \equiv 3 \pmod 4\), let \(f \in \Oc[[X_1, X_2]]\), and let \(D_1\), \(D_2\) be the invariant derivations in the two variables.  Then \((D_1^a D_2^w f)(0) \in p^{\lfloor a/p \rfloor}\Oc\).
\end{problem}

At \(p \equiv 1 \pmod 4\) there is no such gain: \(\lambda_p\) is then a unit, since \(\lambda_p \equiv p\,[X^p]\,\ell = 4^{-m}\binom{2m}{m}\) for \(p = 4m+1\) by the argument of Problem~\ref{cm-pr:lambda}, and for instance \((D^5X)(0) = -12\).  Katz studies congruences for invariant differential operators on Lubin--Tate formal groups of arbitrary finite height \cite[\S II]{Ka}.

\section{The higher odd levels}\label{sec:odd}

At \(p = 3\), the integer \((D^{29}X)(0)\) is divisible exactly by \(3^{10}\), one more than \(\lfloor 29/3 \rfloor = 9\).  Where does the extra factor come from?  Throughout this section \(p \equiv 3 \pmod 4\) and \(q := p^2\), and we write
\[
  E_p(n) := \sum_{\substack{i \ge 1 \\ i \text{ odd}}} \Bigl\lfloor \frac{n}{p^i} \Bigr\rfloor .
\]
The answer is that \(D^n\) gains \(p^{E_p(n)}\).

\begin{problem}\label{pr:log-integral}
\([X^{4m+1}]\, \ell = \binom{2m}{m} / \bigl(4^m (4m+1)\bigr)\), and \(\ell\) has \(p\)-integral coefficients in all degrees less than \(q\).
\end{problem}

\begin{problem}\label{pr:log-q}
\(p\, [X^q]\, \ell \equiv -1 \pmod p\).  (For \(p = 3\): \([X^9]\, \ell = 1/24\).)
\end{problem}

By Problems~\ref{pr:log-integral} and~\ref{pr:log-q}, the truncation \(T\) of \(\ell(X) + X^q/p\) to degree at most \(q\) lies in \(X + X^2\Z_{(p)}[X]\).  It is therefore a parameter on \(\hat E\) over \(\Z_{(p)}\), and \(\Oc[[T]] = \Oc[[X]]\).  In this parameter
\[
  \ell = T - \frac{T^q}{p} + O(T^{q+1}),
  \qquad
  g(z) := \ell^{-1}(z) = z + \frac{z^q}{p} + O(z^{q+1}),
\]
where \(O(\cdot)\) allows arbitrary rational coefficients.  Write the group law in \(T\) again as \(\oplus\), and define operators \(\Delta_j\) on \(\Oc[[T]]\) by \(f(T \oplus Y) = \sum_{j \ge 0} (\Delta_j f)(T)\, Y^j\).  They commute with each other and with \(D\), and \(\Delta_1 = D\).

\begin{problem}\label{pr:Dq}
For \(f \in \Oc[[T]]\), \(f(T \oplus g(z)) = \sum_{n \ge 0} (D^n f)(T)\, z^n/n!\).  Consequently \(D^p = p!\, \Delta_p\) and \(D^q = q!\, (\Delta_q + D/p)\).
\end{problem}

\begin{problem}\label{pr:Dq-unit}
\(\varepsilon := q!/p^{p+1}\) is a \(p\)-adic unit, and \(D^q = p^p \varepsilon\, H\) with \(H := D + p\Delta_q\).
\end{problem}

\begin{problem}\label{pr:pth-powers}
Let \(A\), \(B\) be commuting \(\Oc\)-linear operators on a torsion-free \(\Oc\)-module \(V\) with \((A - B)V \subseteq pV\).  Then \((A^p - B^p)V \subseteq p^2V\).
\end{problem}

Let \(e_j := E_p(p^j) = p^{j-1} + p^{j-3} + \dotsb\), the sum over the exponents \(j - i\) with \(i \le j\) odd.  Thus \(e_0 = 0\) and \(e_1 = 1\).  For \(j \ge 2\) one has \(e_j = p^{j-1} + e_{j-2}\), and for \(j \ge 1\) one has \(e_j \le p\,e_{j-1} + 1\).

\begin{problem}\label{pr:induction}
For every \(j \ge 0\), as operators on \(\Oc[[T]]\),
\[
  D^{p^j} \equiv 0 \pmod{p^{e_j}},
  \qquad
  H^{p^j} \equiv D^{p^j} \pmod{p^{e_j + 1}} .
\]
\end{problem}

\begin{problem}\label{pr:digits}
\(D^n\, \Oc[[X]] \subseteq p^{E_p(n)}\, \Oc[[X]]\) for every \(n \ge 0\).
\end{problem}

Problems~\ref{pr:log-integral}--\ref{pr:digits} prove the following.

\begin{theorem}\label{thm:single}
Let \(p \equiv 3 \pmod 4\) and \(f \in \Oc[[X]]\).  Then \(v_p\bigl((D^n f)(0)\bigr) \ge E_p(n)\) for every \(n \ge 0\).  The same holds for \(f \in \Oc[[X_1, X_2]]\) and \(D = D_1\), the invariant derivation in the first variable.
\end{theorem}

The method gains nothing at the even levels.  In Problem~\ref{pr:Dq-unit}, \(D^q\) has the factor \(p^p\), not the \(p^{p+1}\) one would get from \(\lfloor q/p \rfloor + \lfloor q/q \rfloor\).  The next odd level is \(p^3\).

\section{One theta function}\label{sec:one-theta}

The function \(F\) of Section~\ref{sec:moments} is not a function on \(E\).  Does it nevertheless gain?  At \(p = 3\) the valuations \(v_3(F^{(k)}(0))\) for \(k = 4, 8, \dotsc, 60\) are
\[
  0,\ 0,\ 1,\ 1,\ 2,\ 2,\ 4,\ 4,\ 4,\ 4,\ 4,\ 5,\ 5,\ 7,\ 7,
\]
far below \(\lfloor k/3 \rfloor\).

\begin{problem}\label{pr:logF}
\((\log F)'' = 2\sll^2\).
\end{problem}

Let \(b_n := (D^n X^2)(0) = n!\, [u^n]\, \sll^2 u\), an element of \(\Z[\tfrac12]\).  By Problem~\ref{pr:logF},
\[
  \log F(u) = 2 \sum_{n \ge 0} b_n \frac{u^{n+2}}{(n+2)!} .
\]

\begin{problem}\label{pr:bp}
Let \(p \equiv 3 \pmod 4\) and \(m := (p-3)/4\).  Then \(b_{p-1} \equiv -4^{-m}\binom{2m}{m} \pmod p\).  In particular \(b_{p-1}\) is a \(p\)-adic unit.  (For \(p = 3\): \(b_2 = 2\).)
\end{problem}

\begin{problem}\label{pr:failure}
Let \(p \equiv 3 \pmod 4\).  Then \(F^{(p+1)}(0) \equiv 2 b_{p-1} \pmod p\).  Hence \(F\), written as a series in \(X = \sll u\), does not have \(p\)-integral coefficients.  (For \(p = 3\): \(F^{(4)}(0) = 4\).)
\end{problem}

For \(p \equiv 3 \pmod 4\), the coefficient of \(u^{p+1}\) in \(\log F\) is \(c_p/p\) with \(c_p = 2b_{p-1}/\bigl((p+1)(p-1)!\bigr)\) a unit.  For \(p = 3, 7, 11\) this coefficient is \(\tfrac16, -\tfrac1{140}, \tfrac1{2475}\), so \(c_p = \tfrac12, -\tfrac1{20}, \tfrac1{225}\).  By the Exercise below, for \(g(u) := \exp(c_p u^{p+1}/p)\) one has \(v_p(g^{(k)}(0)) = k/(p^2-1) + O(\log k)\) as \(k\) tends to infinity through multiples of \(p+1\).  At \(p=3\) the comparison gives \(60/8=7.5\), while \(v_3(F^{(60)}(0))=7\).  We do not know whether \(v_p(F^{(k)}(0)) = k/(p^2-1) + O(\log k)\) holds for each fixed prime \(p \equiv 3 \pmod 4\) as \(k\) tends to infinity through multiples of \(4\); the other derivatives vanish at \(0\) by Problem~\ref{pr:F}.

For \(p \equiv 3 \pmod 4\), the first term of the \(u\)-series of \(\log F\) with a non-\(p\)-integral coefficient is \(c_p u^{p+1}/p\).  We next consider the product \(\Theta\) in two variables.  Chapter~\ref{ip-chap}, \emph{The isotropic plane}, shows that for \(p \equiv 3 \pmod 8\) its restrictions to the axes in the directions \(e_\pm\) are integral in the lemniscate parameter (Theorem~\ref{ip-thm:line}).  After division by the two axis restrictions, the remaining factor equals \(1\) on both axes (Problem~\ref{ip-pr:factor}).  At \(p = 7\) the axis restriction is not integral (Problem~\ref{ip-pr:p7}).

\subsection*{Exercise}
Let \(c\) be a \(p\)-adic unit and let \(g(u) := \exp(c u^{p+1}/p)\).  For an integer \(m \ge 0\), write \(k := (p+1)m\).  Then
\[
  v_p\bigl(k!\, [u^k]\, g\bigr) = \frac{k}{p^2 - 1} + \frac{s_p(m) - s_p(k)}{p-1},
\]
where \(s_p\) is the sum of base-\(p\) digits.

\section*{Hints}

\subsection*{\S\ref{sec:moments}}
\begin{description}
\item[\ref{pr:wp-sl}] Both sides are even, have leading term \(u^{-2}\), and satisfy \(f'^2 = 4f^3 - 4f\).
\item[\ref{pr:F}] \(\sigma(\ii u) = \ii\sigma(u)\) and \(\wp(\ii u) = -\wp(u)\).  For the coefficients, \((\log\sigma)'' = -\wp\) and Problem~\ref{pr:wp-sl}.
\item[\ref{pr:gauss}] Factor the Gaussian sum into two one-dimensional theta functions.  Use Problem~\ref{ec-pr:T-fixed} after checking the transformation law.
\item[\ref{pr:moments}] Apply Problem~\ref{pr:gauss} to each coordinate, with \(z_1 = s + t\) and \(z_2 = \zeta(s - t)\).
\item[\ref{pr:t-sym}] \(F(\zeta v) = \sum_k f_{4k} (-1)^k v^{4k}\) if \(F = \sum_k f_{4k} v^{4k}\).  For the congruence, compare \(\Theta(\zeta s e_+ - \zeta t e_-)\) with \(\Theta(s e_+ + t e_-)\).
\end{description}

\subsection*{\S\ref{sec:lemniscate}}
\begin{description}
\item[\ref{pr:sl-series}] Differentiate \(\sll'^2 = 1 - \sll^4\) to get \(\sll'' = -2\sll^3\).
\item[\ref{pr:fg}] \(\sll' u = \sqrt{1 - X^4}\) lies in \(\Z[\tfrac12][[X]]\).
\item[\ref{pr:O-action}] Use the form of \([p]\) modulo \(p\) from \cite[\S IV.7]{Si1}, continuity in the \((p,X)\)-adic topology, and \([\ii]\).
\end{description}

\subsection*{\S\ref{sec:Dp}}
\begin{description}
\item[\ref{pr:Dp-derivation}] Leibniz's rule.
\item[\ref{pr:invariant}] Write \(\delta = g(X)\,d/dX\).  Apply \(\delta\tau_Y = \tau_Y\delta\) to \(X\), with \(\delta\) acting on the \(X\)-variable, and set \(X=0\).
\item[\ref{pr:Dp}] Use Problem~\ref{pr:sl-series}.  Alternatively, compare \(D^p\) with its pullback by \([\ii]\).
\item[\ref{pr:first-gain}] Problem~\ref{pr:Dp} in the variable \(X_1\).
\end{description}

\subsection*{\S\ref{sec:odd}}
\begin{description}
\item[\ref{pr:log-integral}] If \(p \mid 4m+1\) and \(4m + 1 < q\), then \(m \equiv (3p-1)/4 \pmod p\).  Count carries in \(m + m\).
\item[\ref{pr:log-q}] Write \(p = 4r + 3\) and \(m = (q-1)/4 = rp + 3r + 2\).  Take the unit parts of the factorials in \(\binom{2m}{m}\) and use Wilson's theorem.
\item[\ref{pr:Dq}] In the coordinate \(z = \ell(T)\), translation by \(g(z')\) is \(z \mapsto z + z'\).  For the consequences, compare coefficients of \(z^p\) and \(z^q\) in \(\sum_j (\Delta_j f)(T)\, g(z)^j\).
\item[\ref{pr:Dq-unit}] Legendre's formula.
\item[\ref{pr:pth-powers}] Factor \(A^p - B^p\).
\item[\ref{pr:induction}] Induction on \(j\), using Problems~\ref{pr:Dq-unit} and~\ref{pr:pth-powers}; apply the latter to \(H^{p^{j-1}}\) and \(D^{p^{j-1}}\) divided by \(p^{e_{j-1}}\).
\item[\ref{pr:digits}] Write \(n\) in base \(p\).
\end{description}

\subsection*{\S\ref{sec:one-theta}}
\begin{description}
\item[\ref{pr:logF}] \((\log\sigma)'' = -\wp\), \(\wp'' = 6\wp^2 - 2\), \(\wp'^2 = 4\wp^3 - 4\wp\), and Problem~\ref{pr:wp-sl}.
\item[\ref{pr:bp}] By Lagrange inversion, \([u^n]\, \sll^k u = \tfrac{k}{n}\, [X^{n-k}]\, (X/\ell(X))^n\).  Modulo \(p\) and modulo \(X^p\), \((X/\ell(X))^p \equiv 1\).
\item[\ref{pr:failure}] In \(\exp(\log F)\), only the term \(2b_{p-1}u^{p+1}/(p+1)!\) of \(\log F\) has \(p\) in its denominator up to degree \(p + 1\).  For the second statement, compare with Problem~\ref{pr:first-gain}.
\end{description}

\begin{thebibliography}{Si2}
\bibitem[Ka]{Ka} N.~M.~Katz, \emph{Divisibilities, congruences, and Cartier duality}, J.\ Fac.\ Sci.\ Univ.\ Tokyo Sect.\ IA Math.\ \textbf{28} (1981), 667--678.
\bibitem[LT]{LT} J.~Lubin and J.~Tate, \emph{Formal complex multiplication in local fields}, Ann.\ of Math.\ \textbf{81} (1965), 380--387.
\bibitem[Si1]{Si1} J.~H.~Silverman, \emph{The Arithmetic of Elliptic Curves}, 2nd ed., Grad.\ Texts in Math.\ 106, Springer, 2009.
\bibitem[Si2]{Si2} J.~H.~Silverman, \emph{Advanced Topics in the Arithmetic of Elliptic Curves}, Grad.\ Texts in Math.\ 151, Springer, 1994.
\end{thebibliography}

\end{document}
